\section{Conclusion}
\label{sec:conclusion}

We built a verified pipeline from public raw data and a continual benchmark, with the protocol fixed in advance, for forecasting the SOH of second-life lithium-ion cells one to three RPTs ahead. We applied it to six NMC second-life cells and to 20 LFP cells in seven usage conditions.

The main findings are:
\begin{enumerate}
  \item \emph{Learned forecasters clearly beat persistence} on both datasets.
  \item \emph{Forgetting appears only when the tasks differ.} When it does, EWC significantly reduces the error and forgetting of a Temporal Fusion Transformer.
  \item \emph{The proposed TFT\,+\,EWC is not the most accurate forecaster} on either dataset. A usage-only linear model (S1) and fade-rate extrapolation and gradient boosting (S6) are as good or better.
\end{enumerate}

For second-life assessment this means two things. Forecasts should be benchmarked against simple usage- and trend-based references. Continual-learning machinery pays off only on heterogeneous streams, and even there its benefit must be weighed against simpler models that see all data.
